% !TEX root = ../../ctfp-print.tex

\lettrine[lhang=0.17]{程}{序员们围}绕单子发展出了一整套神话. 据说它是编程里最抽象、
最难的概念之一. 有人
``开悟'' 了, 有人没有. 对许多人来说, 理解单子概念的那一刻宛如一场神秘体验. 单子抽象出了如此多样的
构造的精髓, 以至于我们在日常生活里根本找不到一个好的类比. 我们只剩下在黑暗中摸索 --- 就像那些摸象的
盲人, 各自摸到大象的不同部位, 然后得意地喊道: ``是绳子!'' ``是树干!'' 或者 ``是卷饼!''

让我把话说清楚: 围绕单子的种种神秘, 全是一场误解的结果. 单子是一个非常简单的概念. 造成困惑的,
是单子应用的多种多样.

为这个项目做研究时, 我查了查管道胶带 (duct tape, 又写作 duck tape) 和它的种种用途. 这里是
你能拿它做的事情中的一小部分:

\begin{itemize}
  \tightlist
  \item
        密封管道
  \item
        修好阿波罗 13 号飞船上的二氧化碳洗涤器
  \item
        治寻常疣
  \item
        解决苹果 iPhone 4 的掉话问题
  \item
        做一条舞会礼服
  \item
        建造一座吊桥
\end{itemize}

\noindent
现在想象一下, 你并不知道管道胶带是什么, 只想凭这份清单把它琢磨出来. 祝你好运!

所以我想在 ``单子就像\ldots{}'' 这一套陈词滥调里再添一条: 单子就像管道胶带. 它的用途天差地别,
但它的原理非常简单: 它把东西粘在一起. 更确切地说, 它把东西组合起来.

这部分解释了许多程序员 --- 尤其是那些命令式背景出身的人 --- 理解单子的困难. 问题在于, 我们
不习惯用函数组合来思考编程. 这可以理解: 我们常常给中间值起名, 而不是把它们直接从一个函数递给
另一个函数; 我们也倾向于把短短几段胶合代码内联进去, 而不是抽象成辅助函数. 下面是 C 语言里向量
长度函数的命令式风格实现:

\begin{snip}{cpp}
double vlen(double * v) {
    double d = 0.0;
    int n;
    for (n = 0; n < 3; ++n)
        d += v[n] * v[n];
    return sqrt(d);
}
\end{snip}
把它和 (风格化的) Haskell 版本比一比, 后者把函数组合明确地写了出来:

\src{snippet01}
(这里为了让事情显得更加晦涩, 我把幂运算符 \code{(\^{})} 的第二个参数设成 \code{2}, 做了
部分应用.)

我并不是说 Haskell 的无点风格总是更好, 我只是说函数组合是我们编程中所做一切的根本. 而尽管
我们实际上一直在组合函数, Haskell 仍然费了大力气去提供一种命令式风格的语法 --- \code{do}
记法 --- 用来做单子组合. 我们稍后会看到它的用处. 但首先, 让我解释为什么我们本来就需要单子式的
组合.

\section{Kleisli 范畴}

我们先前是通过装饰普通函数得到 \hyperref[kleisli-categories]{写者单子}的. 那种特定的装饰做法,
是把它们的返回值与字符串配对, 或者更一般地, 与某个幺半群的元素配对. 我们现在可以认出, 这样的装饰
其实是一个函子:

\src{snippet02}
随后我们找到了组合装饰过的函数 (即 Kleisli 箭头) 的办法, 它们是形如下面的函数:

\src{snippet03}
正是在这个组合里面, 我们实现了日志的累积.

现在我们可以给出 Kleisli 范畴更一般的定义了. 我们从范畴 $\cat{C}$ 和自函子 $m$ 出发.
相应的 Kleisli 范畴 $\cat{K}$ 与 $\cat{C}$ 有着同样的物件, 但态射不同. $\cat{K}$ 中两个
物件 $a$ 和 $b$ 之间的态射, 是在原范畴 $\cat{C}$ 中实现的这样一个态射:
\[a \to m\ b\]
要牢牢记住: 我们把 $\cat{K}$ 中的一个 Kleisli 箭头当作 $a$ 与 $b$ 之间的态射, 而不是
$a$ 与 $m\ b$ 之间的态射.

在我们的例子里, $m$ 特化为 \code{Writer w}, 其中 \code{w} 是某个固定的幺半群.

只有当我们能为 Kleisli 箭头定义恰当的组合时, 它们才构成一个范畴. 如果存在这样的组合 ---
它满足结合律, 并且为每个物件提供恒等箭头 --- 那么函子 $m$ 就叫作一个 \newterm{monad 单子},
由此得到的范畴就叫作 Kleisli 范畴.

在 Haskell 里, Kleisli 组合用鱼运算符 \code{>=>} 来定义, 恒等箭头则是一个叫 \code{return}
的多态函数. 下面就是用 Kleisli 组合给出的单子定义:

\src{snippet04}
请记住, 定义单子的等价方式有很多种, 而这并不是 Haskell 生态里最主要的那一种. 我喜欢它, 是因为
它在概念上简单、能给出直觉; 但另有的定义在实际编程时更方便. 我们马上就会讲到它们.

在这个表述里, 单子律非常容易表达. Haskell 无法强制它们成立, 但我们可以拿它们做等式推理.
它们无非就是 Kleisli 范畴的标准组合律:

\begin{snip}{haskell}
(f >=> g) >=> h = f >=> (g >=> h) -- associativity
return >=> f = f                  -- left unit
f >=> return = f                  -- right unit
\end{snip}
这类定义也说清了单子究竟是什么: 它是一种组合装饰过的函数的办法. 它无关副作用或状态, 它关乎的
是组合. 我们稍后会看到, 装饰过的函数可以用来表达各种副作用或状态, 但那不是单子之所以为单子的地方.
单子就是那条粘糊糊的管道胶带, 把一个装饰函数的一端粘到另一个装饰函数的一端上.

回到我们的 \code{Writer} 例子: 记日志的那些函数 (\code{Writer} 函子的 Kleisli 箭头)
构成一个范畴, 因为 \code{Writer} 是一个单子:

\src{snippet05}
只要 \code{w} 的幺半群律成立, \code{Writer w} 的单子律就成立 (幺半群律在 Haskell 里同样
没法强制).

\code{Writer} 单子还定义了一个很有用的 Kleisli 箭头, 叫 \code{tell}. 它唯一的用途就是把
自己的参数加进日志:

\src{snippet06}
稍后我们会把它当成积木, 用来搭其他单子函数.

\section{鱼的解剖}

在给不同的单子实现鱼运算符时, 你很快就会发觉有大量代码在重复, 而且很容易就能抽出去. 首先,
两个函数的 Kleisli 组合必须返回一个函数, 所以它的实现不妨从一个接受 \code{a} 类型参数的 lambda
开始:

\src{snippet07}
我们对这个参数唯一能做的事, 就是把它传给 \code{f}:

\src{snippet08}
到这一步, 我们得产出类型为 \code{m c} 的结果, 而手头只有一个类型为 \code{m b} 的物件和一个
函数 \code{g :: b -> m c}. 我们来定义一个替我们干这活的函数. 这个函数叫作 \emph{bind} (绑定),
通常写成中缀运算符的形式:

\src{snippet09}
对每个单子来说, 我们不必定义鱼运算符, 而可以改去定义 \code{bind}. 事实上, Haskell 里单子的
标准定义用的就是 \code{bind}:

\src{snippet10}
这里是 \code{Writer} 单子的 \code{bind} 定义:

\src{snippet11}
它确实比鱼运算符的定义更短.

我们还可以进一步剖析 \code{bind}, 利用 \code{m} 是一个函子这一事实. 我们可以用
\code{fmap} 把函数 \code{a -> m b} 施加到 \code{m a} 的内容上. 这会把 \code{a} 变成
\code{m b}. 于是应用的结果类型为 \code{m (m b)}. 这不正好是我们想要的 --- 我们需要的结果类型是
\code{m b} --- 但已经很接近了. 我们缺的只是一个函数, 把 \code{m} 的双重应用折叠、压平. 这样的
函数叫作 \code{join}:

\src{snippet12}
有了 \code{join}, 我们可以把 \code{bind} 重写为:

\src{snippet13}
这就把我们引向定义单子的第三种选择:

\src{snippet14}
这里我们明确地要求了 \code{m} 是一个 \code{Functor}. 在前两个单子的定义里我们并不需要这么
做. 那是因为任何支持鱼运算符或 \code{bind} 运算符的类型构造子 \code{m} 都自动是一个函子. 举例
来说, 用 \code{bind} 和 \code{return} 就能定义出 \code{fmap}:

\src{snippet15}
为了完整, 这里给出 \code{Writer} 单子的 \code{join}:

\src{snippet16}

\section{\texttt{do} 记法}

写单子代码的一种方式是与 Kleisli 箭头打交道 --- 用鱼运算符把它们组合起来. 这种编程模式是无点
风格的推广. 无点代码紧凑, 而且常常相当优雅. 但一般而言, 它可能难以理解, 近乎晦涩. 这就是为什么
大多数程序员更喜欢给函数参数和中间值起名.

处理单子时, 这就意味着偏爱 \code{bind} 运算符胜过鱼运算符. \code{bind} 取一个单子值, 返回
一个单子值. 程序员可以选择给这些值起名字. 但这几乎算不上什么改进. 我们真正想要的, 是假装自己在跟
普通值打交道, 而不是跟封装着它们的单子容器打交道. 命令式代码就是这么工作的 --- 副作用, 比如更新
一个全局日志, 大体上藏在视线之外. 而这正是 \code{do} 记法在 Haskell 里模拟出来的东西.

那么你可能会问, 干吗还要用单子? 如果我们想让副作用隐形, 何不干脆留在命令式语言里? 答案是:
单子给了我们对副作用好得多的控制. 举例来说, \code{Writer} 单子里的日志是在函数之间传递的,
从不全局暴露. 根本没有把日志搞乱、或制造出数据竞争的可能. 而且单子代码界限分明, 与程序的其余部分
隔离开来.

\code{do} 记法不过是单子组合的语法糖. 表面上它很像命令式代码, 但它直接翻译成一系列 \code{bind}
与 lambda 表达式.

举例来说, 取我们先前用来阐释 \code{Writer} 单子里 Kleisli 箭头组合的那个例子. 用我们现在的
定义, 它可以重写为:

\src{snippet17}
这个函数把输入字符串里的所有字符转成大写, 并把它切成单词, 全程还产出一份记录自身行为的日志.

在 \code{do} 记法下, 它看起来会是这样:

\src{snippet18}
这里 \code{upStr} 只是一个 \code{String}, 尽管 \code{upCase} 产出的是一个 \code{Writer}:

\src{snippet19}
这是因为 \code{do} 块会被编译器脱糖 (desugar) 成:

\src{snippet20}
\code{upCase} 的单子结果被绑定到一个接受 \code{String} 的 lambda 上. 出现在 \code{do}
块里的, 正是这个字符串的名字. 当我们读到这一行:

\src{snippet21}
我们就说 \code{upStr} \emph{得到} 了 \code{upCase s} 的结果.

伪命令式的风格在我们内联 \code{toWords} 时表现得更加明显. 我们把它换成对 \code{tell} 的调用
--- 它把字符串 \code{"toWords "} 记进日志 --- 紧接着是对 \code{return} 的调用, 带着用
\code{words} 切分字符串 \code{upStr} 得到的结果. 注意 \code{words} 是一个在字符串上工作的
普通函数.

\src{snippet22}
在这里, \code{do} 块中的每一行都会在脱糖后的代码里引入一层新的嵌套 \code{bind}:

\src{snippet23}
注意 \code{tell} 产出的是一个单位值, 所以它不必被传给后面的那个 lambda. 忽略一个单子结果的
内容 (但不忽略它的效应 --- 这里是对日志的贡献) 相当常见, 于是有一个专门的运算符, 在这种情形下取代
\code{bind}:

\src{snippet24}
我们这段代码真正的脱糖结果是这样:

\src{snippet25}
一般来说, \code{do} 块由若干行 (或子块) 组成: 它们要么用左箭头引入新的名字, 供代码其余部分
使用, 要么纯粹为了副作用而执行. \code{bind} 运算符隐含在各行代码之间. 顺便说一句, 在 Haskell
里可以把 \code{do} 块的排版换成花括号和分号. 这正是 ``把单子说成是重载分号的办法'' 这种说法的
依据.

注意, 在脱糖 \code{do} 记法时 lambda 与 \code{bind} 运算符的嵌套, 其效果是让每一行的结果去
影响 \code{do} 块其余部分的执行. 这个性质可以用来引入复杂的控制结构, 比如说模拟异常.

有趣的是, \code{do} 记法的等价物已经在命令式语言里找到了用武之地, C++ 尤其如此. 我说的是
可恢复函数 (resumable functions) 或协程 (coroutines). 这并不是什么秘密: C++ 的
\urlref{https://bartoszmilewski.com/2014/02/26/c17-i-see-a-monad-in-your-future/}
{future 构成一个单子}. 它是续体单子的一个例子, 我们很快就会
讨论. 续体的麻烦在于
它们非常难以组合. 在 Haskell 里, 我们用 \code{do} 记法把 ``我的处理器会去调用你的处理器''
这团面条, 变成看上去极为接近顺序代码的东西. 可恢复函数在 C++ 里让同样的变换成为可能. 同一套机制还
可以用来把
\urlref{https://bartoszmilewski.com/2014/04/21/getting-lazy-with-c/}{嵌套循环这团面条}
变成列表推导式或 ``生成器'', 它们本质上就是
列表单子的 \code{do} 记法. 若没有单子这层统一的抽象, 这些问题里的每一个通常
都要靠给语言提供自定义扩展来
解决. 而在 Haskell 里, 这一切都交给库来办.
